
We need a bit more terminology. Let \Rel be the category of sets and binary relations. Every partial function can be considered as a binary relation and the inclusion of $\Par$ in $\Rel$ is functorial, we shall denote 
by $U$ as it 'forgets' the many-one nature of a partial function so
that
$\Par \xrightarrow{U} \Rel$ which. 

If $F: \catc \morph \Par$ is a functor then composition  yields a functor into $\Rel$ which we shall denote $F_R$ so that
 $F_R : \catc \morph \Rel$.

 I shall use $R^{-1}$ for the converse of a relation. I require some simple facts.

 \begin{lemma}
 \llabel{threesidedlemma}
 If in the category \Par the following diagram of partial functions commutes 
$
\begin{array}{c p{0.25cm} c  p{0.25cm} c }
             &&   \Rnode{c}{c} &&              \\[0.4cm]
\Rnode{b1}{b_1} &&                   && \Rnode{b2}{b_2} \\[0.4cm]
             &&   \Rnode{a}{a} &&              
\end{array} 
\begin{arrows}
\ncarr{a}{b1}
\alabel{f_1}
\ncarr{a}{b2}
\blabel{f_2}
\ncarr{b1}{c}
\alabel{g_1} 
\ncarr{b2}{c}
\blabel{g_2}
\end{arrows}
$
commutes then  the diagram of relationships
$
\begin{array}{c p{0.25cm} c  p{0.25cm} c }
             &&   \Rnode{c}{c} &&              \\[0.4cm]
\Rnode{b1}{b_1} &&                   && \Rnode{b2}{b_2} \\[0.4cm]
             &&   \Rnode{a}{a} &&              
\end{array} 
\begin{arrows}
\ncarr{b1}{a}
\blabel{U(f_1)^{-1}}
\ncarr{a}{b2}
\blabel{U(f_2)}
\ncarr{b1}{c}
\alabel{U(g_1)} 
\ncarr{b2}{c}
\blabel{U(g_2)}
\end{arrows}
$
commutes in the category \Rel.
 \end{lemma}

 \begin{lemma}
 \llabel{atleasttwopartialfunctionslemma}

 If $R: x \morph y$ in the category of sets 
 such that $R \circ U(q)=U(f)$ and if $R$ is not single-valued
 then there exists at least  two distinct partial functions $r: x \morph y$
 such that $r \circ q = f$.
 \end{lemma}

 \begin{lemma}
 \llabel{whatsit}
 Suppose I have the following 'zigzag' of partial functor in \Par:
\workt{I dont have macros for zigzags the right way around so this is an approximation}
$$
\arraycolsep=7pt
\begin{array}{cccccc}
 &\bOb{x}& &\bOb{z} && upto fn\\[0.8cm]
\bOb{w}& &\bOb{y}&
\end{array}
\begin{arrows}
\ncarr{w}{x}\alabel{f_0}[0.5][0.4]
\ncarr{y}{x}\blabel{f'_0}[0.5][0.4]
\ncarr{y}{z}\blabel{f_1}[0.5][0.4]
\end{arrows}
$$
\newt{if this zigzag is tight}
then if the relation 
$U(f_0)\circ U(f'_0)^{-1} \circ U(f_1) \circ U(f'{_1})^-1... \circ U(f_n)$
is single-valued (it is in the image of $U$ therefore) 
then for each $i$, $1 \leq i \leq n-1$, the relation 
$U(f'_i) \circ U(f'_{i+1})$ is single-valued.
\workt{I am going to have to move this lemma to the end of sticklebacks I think so that I don't have to explain tightness here.}
 \end{lemma}
